\subsection{四元数体}

由定理 1 推出，域 R 是一个极大有序域，因而（在同构意义下）R 上有限秩的唯一非交换除环就是 R 上的四元数体；它记作 H，称为实四元数体（或在无混淆之虞时简称四元数体）。由于 H 在域 R 上的秩为 4，我们可以在 H 上定义一个与 \(R^{4}\) 的拓扑同胚的拓扑（第六章，§ 1，第 5 号）。确切地说，我们通常把 H 与 \(R^{4}\) 等同，即把 H 的标准基的元素 1、i、j、k 分别等同于向量 \(e_{0}, e_{1}, e_{2}, e_{3}\) （\(R^{4}\) 的标准基的向量）。

我们回顾，H 的标准基的乘法表由下列公式给出

\[
\begin{array}{c} i ^ {2} = j ^ {2} = k ^ {2} = - \mathrm{I}, \quad i j = - j i = k, \\ j k = - k j = i, \quad k i = - i k = j. \end{array}
\]

H 的拓扑不仅与 H 的环结构相容（第六章，§ 1，第 5 号），而且与它的除环结构相容；因为若 x 是一个非零四元数，则 \(x^{-1}\) 的坐标是 x 的坐标的有理函数，其分母不为零。因此，赋有这一拓扑的除环 H 是一个非交换的拓扑除环。四元数 \(a + bi\) （a、b 为实数）构成 H 的一个（拓扑）子域，它同构于域 C，并常与 C 等同。

这样我们就有了局部紧、连通的拓扑除环的第三个例子，其余两个是 R 与 C。事实上，具有这两条性质的拓扑除环仅此三个。

我们从代数知道，四元数 \(x = x_{0} + x_{1}i + x_{2}j + x_{3}k\) 的约化范数是

\[
\mathrm{N} (\boldsymbol {x}) = x _ {0} ^ {2} + x _ {1} ^ {2} + x _ {2} ^ {2} + x _ {3} ^ {2} = | | \boldsymbol {x} | | ^ {2}
\]

（因此它是 x 的欧氏范数的平方）。由于

\[
\mathrm{N} (\boldsymbol {x y}) = \mathrm{N} (\boldsymbol {x}) \mathrm{N} (\boldsymbol {y}),
\]

由此得出，范数为 1 的所有四元数所成的集——它与球面 \(S_{3}\) 相同——构成非零四元数的乘法群 \(H^{*}\) 的一个紧子群。

命题 2. 非零四元数的乘法群 \(\mathbf{H}^*\) 是一个拓扑群，同构于它的子群 \(\mathbf{R}_+^*\) 与 \(\mathbf{S}_3\) 的乘积。

每个四元数 \(x \neq 0\) 都可以写成 \(x = ||x||.z\) ，其中 \(z\) 是范数为 1 的四元数；由于 \(||xx'|| = ||x||.||x'||\) ，映射 \(x \to (||x||, x/||x||)\) （\(H^*\) 到 \(R_+^* \times S_3\) 上的）是群结构的一个同构，并且由第六章 § 2 第 3 号命题 3，它是 \(H^*\) 到 \(R_+^* \times S_3\) 上的一个同胚。

注。1) 利用关系式 \(||x + y|| \leqslant ||x|| + ||y||\) 与 \(||xy|| = ||x||.||y||\) ，可以像第 2 号中对复数域那样直接证明 \(\mathbf{R}^4\) 的拓扑与 H 的除环结构相容（参看第九章，§ 3，第 2 号）。

\begin{enumerate}
\def\labelenumi{\arabic{enumi})}
\setcounter{enumi}{1}
\item
  由已证的结果推出，球面 \(S_{1}\) 与 \(S_{3}\) 能够带有与它们的拓扑相容的群结构。我们后面将看到，对每个异于 1 和 3 的整数 n，\(S_{n}\) 上不存在与 \(S_{n}\) 的拓扑相容的群结构。
\item
  群 \(\mathbf{S}_3\) 的每个点都有一个同胚于 \(\mathbf{R}^3\) 的邻域（第六章，§ 2，第 4 号，命题 5），但 \(\mathbf{S}_3\) 不局部同构于 \(\mathbf{R}^3\) ；因为若是这样，则由于它是连通的（第七章，§ 2，第 2 号，定理 1），它将是交换的，而事实并非如此，因为 \(i\) 与 \(j\) 都属于 \(\mathbf{S}_3\) 且 \(ij \neq ji\) （参看第五章，§ 3）。
\end{enumerate}

